% from ega2-1-fr.tex
\oldpage[II]{69}

\subsection{Critères d'immersion dans un spectre homogène.}
\label{subsection:II.3.8-fr}

\begin{env}[3.8.1]
\label{II.3.8.1-fr}
Avec les notations de (\hyperref[II.3.7.1-fr]{3.7.1}), la question de savoir
si $r_{\mathcal{L},\psi}$ est une immersion (resp. une immersion ouverte, une
immersion fermée) est évidemment \emph{locale sur $Y$}.
\end{env}

\begin{proposition}[3.8.2]
\label{II.3.8.2-fr}
Sous les hypothèses de (\hyperref[II.3.7.3-fr]{3.7.3}), pour que
$r_{\mathcal{L},\psi}$ soit partout définie et soit une immersion, il faut et
il suffit qu'il existe une famille de sections $s_\alpha\in S_{n_\alpha}$
($n_\alpha>0$) telle que, si on pose
$f_\alpha=\psi^\flat(s_\alpha)$, les conditions suivantes soient vérifiées~:
\begin{enumerate}
  \item[(i)] Les $X_{f_\alpha}$ forment un recouvrement de $X$.
  \item[(ii)] Les $X_{f_\alpha}$ sont des ouverts affines.
  \item[(iii)] Pour tout $\alpha$ et tout
    $t\in\Gamma(X_{f_\alpha},\mathcal{O}_X)$, il existe un entier $m>0$ et
    un $s\in S_{mn_\alpha}$ tels que
    $t=(\psi^\flat(s)|X_{f_\alpha})(f_\alpha|X_{f_\alpha})^{-m}$.
\end{enumerate}

Pour que $r_{\mathcal{L},\psi}$ soit partout définie et soit une immersion
ouverte, il faut et il suffit qu'il existe une famille $(s_\alpha)$ vérifiant
les conditions (i), (ii), (iii) et~:
\begin{enumerate}
  \item[(iv)] Pour tout $n>0$ et tout $s\in S_{nn_\alpha}$ tel que
    $\psi^\flat(s)|X_{f_\alpha}=0$, il existe un entier $k>0$ tel que
    $s_\alpha^ks=0$.
\end{enumerate}

Pour que $r_{\mathcal{L},\psi}$ soit partout définie et soit une immersion
fermée, il faut et il suffit qu'il existe une famille $(s_\alpha)$ vérifiant
(i), (ii), (iii) et~:
\begin{enumerate}
  \item[(v)] Les $D_+(s_\alpha)$ forment un recouvrement de
    $P=\operatorname{Proj}(S)$.
\end{enumerate}
\end{proposition}

\begin{proof}
En effet, pour que $r$ soit une immersion (resp. une immersion fermée), il
faut et il suffit qu'il existe un recouvrement de $r(G(\psi))$ (resp. de
$P$) par des $D_+(s_\alpha)$ tel que si on pose
$V_\alpha=r^{-1}(D_+(s_\alpha))$, la restriction de $r$ à $U_\alpha$ soit
une immersion fermée de $V_\alpha$ dans $D_+(s_\alpha)$
(\hyperref[I.4.2.4-fr]{I, 4.2.4}). La condition (i) exprime à la fois que
$r$ est partout définie et que les $D_+(s_\alpha)$ recouvrent $r(X)$,
d'après (\hyperref[II.3.7.3.1-fr]{3.7.3.1})~; comme $D_+(s_\alpha)$ est
affine, les conditions (ii) et (iii) expriment que la restriction de $r$ à
$X_{f_\alpha}$ est une immersion fermée dans $D_+(s_\alpha)$, d'après
(\hyperref[I.4.2.3-fr]{I, 4.2.3})~; enfin, comme (iii) et (iv) expriment que
$\psi_{(s_\alpha)}^\flat$ est bijective (notation de
(\hyperref[II.3.7.3.2-fr]{3.7.3.2})), (ii), (iii) et (iv) expriment que la
restriction de $r$ à $X_{f_\alpha}$ est un isomorphisme sur
$D_+(s_\alpha)$ pour tout $\alpha$, donc (i), (ii), (iii) et (iv) expriment
que $r$ est une immersion ouverte.
\end{proof}

\begin{corollary}[3.8.3]
\label{II.3.8.3-fr}
Sous les hypothèses de (\hyperref[II.3.7.6-fr]{3.7.6}), si
$r_{\mathcal{L},\psi'}$ est partout définie et est une immersion, il en est
de même de $r_{\mathcal{L},\psi}$. Si on suppose en outre que $u$ est
(\textbf{TN})-surjectif et si $r_{\mathcal{L},\psi'}$ est une immersion
ouverte (resp. fermée), il en est de même de $r_{\mathcal{L},\psi}$.
\end{corollary}

\begin{proof}
En effet, d'après (\hyperref[II.3.8.2-fr]{3.8.2}), il y a une famille de
$s'_\alpha\in S'_{n_\alpha}$ telle que si on pose
$f_\alpha=\psi'^\flat(s'_\alpha)$, les conditions (i), (ii), (iii) soient
vérifiées. Or, si on pose $s_\alpha=u(s'_\alpha)$, on a aussi
$f_\alpha=\psi^\flat(s_\alpha)$, et si on a
$t=(\psi'^\flat(s')|X_{f_\alpha})(f_\alpha|X_{f_\alpha})^{-m}$, on a aussi
$t=(\psi^\flat(s)|X_{f_\alpha})(f_\alpha|X_{f_\alpha})^{-m}$ en posant
$s=u(s')$, d'où la première assertion. La seconde résulte aussitôt de ce que
$\operatorname{Proj}(u)$ est une immersion fermée.
\end{proof}

\begin{proposition}[3.8.4]
\label{II.3.8.4-fr}
Supposons vérifiées les hypothèses de
(\hyperref[II.3.7.7-fr]{3.7.7}) et en outre que $q:X\to Y$ soit un morphisme
de type fini. Alors, pour que $r_{\mathcal{L},\psi}$ soit partout définie et
soit une immersion, il faut et il suffit qu'il existe $\lambda$ tel que
$r_{\mathcal{L},\psi_\lambda}$ soit partout définie et soit une immersion et
dans ce cas $r_{\mathcal{L},\psi_\mu}$ est partout définie et est une
immersion pour $\mu\geq\lambda$.
\end{proposition}

